\subsection{Практические задания: таймеры и алгоритмы}
Этот блок следует выполнять после разделов про таймеры, неблокирующее выполнение и конечные автоматы. Во всех задачах ниже допустимо использовать \verb|Timer.new(...)|, \verb|Timer.callLater(...)| и глобальные переменные состояния.

\begin{taskbox}[Отложенное включение синего цвета]
\indent Необходимо реализовать алгоритм, при котором светодиодная линейка остаётся выключенной в течение первых трёх секунд работы программы, после чего все светодиоды одновременно загораются синим цветом \verb|{0,0,1}|. Реализация должна использовать механизм отложенного вызова без применения блокирующих задержек.

\noindent\textbf{Ожидаемый результат:} через 3 секунды после запуска программы линейка полностью загорается синим цветом.
\end{taskbox}

\begin{taskbox}[Периодическое мигание красным цветом]
\indent Требуется реализовать периодическое переключение всей светодиодной линейки между состояниями «включено красным цветом \verb|{1,0,0}|» и «выключено \verb|{0,0,0}|» с периодом 0{,}5~с. Алгоритм должен быть основан на использовании асинхронного таймера и логической переменной состояния.

\noindent\textbf{Ожидаемый результат:} равномерное мигание линейки с периодом 0{,}5~с.
\end{taskbox}

\begin{taskbox}[Эффект «бегущего огня» (один активный светодиод)]
\indent Необходимо реализовать алгоритм, при котором в каждый момент времени включён только один светодиод зелёного цвета \verb|{0,1,0}|, а остальные выключены. Активный светодиод должен последовательно перемещаться от индекса \verb|0| к \verb|24| с циклическим повторением.

\noindent\textbf{Ожидаемый результат:} визуальный эффект перемещающегося «огонька» по всей линейке.
\end{taskbox}

\begin{taskbox}[Модель светофора на основе конечного автомата]
\indent Требуется реализовать конечный автомат с тремя состояниями: красный \verb|{1,0,0}|, жёлтый \verb|{1,1,0}| и зелёный \verb|{0,1,0}|. Состояния должны циклически сменять друг друга с заданным временным интервалом.

\noindent\textbf{Ожидаемый результат:} периодическая и стабильная смена цветов по схеме «красный -- жёлтый -- зелёный».
\end{taskbox}

\begin{taskbox}[Плавная пульсация яркости красного канала (PWM-модуляция)]
\indent Необходимо реализовать алгоритм плавного увеличения и уменьшения яркости красного канала от значения \verb|0.0| до \verb|1.0| и обратно. Изменение яркости должно происходить малыми приращениями с использованием таймера.

\noindent\textbf{Ожидаемый результат:} вся светодиодная линейка плавно пульсирует красным цветом.
\end{taskbox}

\begin{taskbox}[Стробоскопический режим (белые вспышки)]
\indent Требуется реализовать быстрое периодическое включение и выключение всей светодиодной линейки белым цветом \verb|{1,1,1}| с малым временным интервалом.

\noindent\textbf{Ожидаемый результат:} частые и чёткие белые вспышки, создающие эффект стробоскопа.
\end{taskbox}

\begin{taskbox}[Инвертируемая полицейская мигалка]
\indent Требуется реализовать алгоритм, при котором цвета чётных и нечётных светодиодов периодически меняются местами с заданным интервалом времени.

\noindent\textbf{Ожидаемый результат:} регулярное чередование цветовых паттернов, имитирующее работу полицейской световой сигнализации.
\end{taskbox}

\begin{taskbox}[Генерация цветового шума]
\indent Необходимо реализовать алгоритм, который с заданной периодичностью устанавливает для каждого светодиода случайные значения компонент RGB в диапазоне от 0 до 1.

\noindent\textbf{Ожидаемый результат:} динамично изменяющаяся цветовая картина без видимой закономерности.
\end{taskbox}

\begin{taskbox}[Индикатор загрузки (loading bar)]
\indent Требуется реализовать алгоритм постепенного заполнения светодиодной линейки от младших индексов к старшим, при котором включённые светодиоды сохраняют своё состояние. Для анимации используйте таймер, который на каждом тике увеличивает индекс последнего включённого светодиода.

\noindent\textbf{Ожидаемый результат:} визуальный эффект «заполняющейся» полосы.
\end{taskbox}

\begin{taskbox}[Сигнал SOS по азбуке Морзе]
\indent Необходимо реализовать световую передачу сигнала «SOS» в соответствии с азбукой Морзе: три коротких импульса, три длинных импульса и снова три коротких. Длительности импульсов должны соответствовать заданным временным параметрам.

\noindent\textbf{Ожидаемый результат:} корректно воспроизводимый световой сигнал SOS.
\end{taskbox}

\begin{taskbox}[Двоичный счётчик секунд]
\indent Требуется реализовать алгоритм отображения текущего значения секунд в двоичном виде, используя светодиоды линейки как индикаторы отдельных битов. Обновление состояния должно происходить один раз в секунду.

\noindent\textbf{Ожидаемый результат:} на линейке отображается двоичное представление времени в секундах.
\end{taskbox}
